% Quelle: 3. Übung Lösungen.pdf
% Art: Übung mit Lösungen
% Themen: Lemma von Bezout, Chinesischer Restsatz, Euklidischer Algorithmus, Eulersche phi-Funktion, Affine Ebene AG(2,p), Bijektionen, Gitterwege, Binomischer Lehrsatz, Vandermonde-Identität
\section{Übung 3 (mit Lösungen)}
\subsection{Lemma von Bezout, Chinesischer Restsatz}

\begin{aufgabe}[Aufgabe 1]
Lösen Sie diese Systeme simultaner Kongruenzen
\begin{enumerate}[label=\alph*)]
\item $z \equiv 3 \pmod{5}$, \quad $z \equiv 4 \pmod{7}$
\item $z \equiv -1 \pmod{12}$, \quad $z \equiv 2 \pmod{5}$
\end{enumerate}
\end{aufgabe}

\begin{loesung}
a) EA für 7, 5 \qquad $z = 5x + 3 = 7y + 4$
\begin{align*}
7 &= 1 \cdot 5 + 2\\
5 &= 2 \cdot 2 + 1\\
2 &= 2 \cdot 1 + 0
\end{align*}
% KORRIGIERT: "Lösingsvarianten"; Hinweis auf ggT ergänzt
Der letzte Rest $\neq 0$ ist $r_3 = 1$, also $\ggT(7, 5) = 1$ und das System ist lösbar. Jetzt gibt es zwei Lösungsvarianten.

Variante 1: EA rückwärts
\begin{align*}
\Rightarrow\quad 1 &= 5 - 2 \cdot 2\\
&= 5 - 2 \cdot (7 - 1 \cdot 5)\\
% KORRIGIERT: "= 18 (mod 35)" statt Kongruenz; "= 4 - 3" als b - a gekennzeichnet
&= 3 \cdot 5 - 2 \cdot 7 = 4 - 3 \quad (= b - a)\\
\Rightarrow\quad z &= 3 \cdot 5 + 3 = 2 \cdot 7 + 4 = 18, \quad \text{also } z \equiv 18 \pmod{35}
\end{align*}
Variante 2: mittels der Matrix $Q$
\[
\begin{pmatrix} m \\ n \end{pmatrix} = Q \begin{pmatrix} r_3 \\ 0 \end{pmatrix}, \quad
Q = \begin{pmatrix} 1 & 1 \\ 1 & 0 \end{pmatrix}\begin{pmatrix} 2 & 1 \\ 1 & 0 \end{pmatrix}\begin{pmatrix} 2 & 1 \\ 1 & 0 \end{pmatrix}
= \begin{pmatrix} 3 & 1 \\ 2 & 1 \end{pmatrix}\begin{pmatrix} 2 & 1 \\ 1 & 0 \end{pmatrix}
= \begin{pmatrix} 7 & 3 \\ 5 & 2 \end{pmatrix}
\]
\[
\operatorname{Det} Q = 7 \cdot 2 - 5 \cdot 3 = -1 = a - b,
\]
\[
% KORRIGIERT: "= 18 (mod 35)" statt Kongruenz
\Rightarrow\quad z = 7 \cdot 2 + 4 = 5 \cdot 3 + 3 = 18, \quad \text{also } z \equiv 18 \pmod{35}
\]

b) EA für 12, 5
\begin{align*}
12 &= 2 \cdot 5 + 2\\
5 &= 2 \cdot 2 + 1
\end{align*}
% KORRIGIERT: Begründung der Lösbarkeit ergänzt
Also $\ggT(12, 5) = 1$, das System ist lösbar.

EA rückwärts \qquad $b - a = 2 - (-1) = 3$
\begin{align*}
\Rightarrow\quad 1 &= 5 - 2 \cdot 2 = 5 - 2 \cdot (12 - 2 \cdot 5)\\
&= 5 \cdot 5 - 2 \cdot 12\\
2 - (-1) &= 5 \cdot 5 \cdot 3 - 2 \cdot 12 \cdot 3\\
% KORRIGIERT: "= -73 (mod 60)" statt Kongruenz
\Rightarrow\quad z &= 2 - 5 \cdot 5 \cdot 3 = (-1) - 2 \cdot 12 \cdot 3 = -73, \quad \text{also } z \equiv -73 \equiv 47 \pmod{60}
\end{align*}
\end{loesung}

\subsection{Eulers $\varphi$-Funktion}

\begin{aufgabe}[Aufgabe 2]
Bestimmen Sie $\varphi(n)$ für
\[
n = 1000, \qquad n = 1001, \qquad n = 1002
\]
\end{aufgabe}

\begin{loesung}
\begin{align*}
n &= 1000 = 2^3 \cdot 5^3\\
% KORRIGIERT: Voraussetzung ggT = 1 für die Multiplikativität ergänzt
\varphi(1000) &= \varphi(2^3) \cdot \varphi(5^3) = (2^3 - 2^2) \cdot (5^3 - 5^2) = 4 \cdot 100 = 400 \quad (\text{da } \ggT(2^3, 5^3) = 1)
\end{align*}
\begin{align*}
n &= 1001 = 7 \cdot 11 \cdot 13\\
% KORRIGIERT: "phi(1000)" statt phi(1001) (vorher als UNSICHER markiert)
\varphi(1001) &= \varphi(7) \cdot \varphi(11) \cdot \varphi(13) = 6 \cdot 10 \cdot 12 = 720 \quad (\text{Faktoren paarweise teilerfremd})
\end{align*}
\begin{align*}
n &= 1002 = 2 \cdot 3 \cdot 167\\
\varphi(1002) &= \varphi(2) \cdot \varphi(3) \cdot \varphi(167) = 1 \cdot 2 \cdot 166 = 332 \quad (\text{Faktoren paarweise teilerfremd})
\end{align*}
\end{loesung}

\subsection{Affine Ebene AG$(2, p)$}

\begin{aufgabe}[Aufgabe 3]
Die affine Ebene AG$(2,5)$ der Ordnung 5 besteht aus den 25 Punkten
\[
\mathrm{AG}(2,5) = \{\, (x, y) \mid x, y \in \Zm{5} \,\}
\]
\begin{enumerate}[label=\alph*)]
\item Welche Punkte bilden die Gerade $y = 3x + 2$?
\item Berechnen Sie den Schnittpunkt $S$ der Geraden $y = 3x + 2$, \ $y = 4x + 1$?
\item In AG$(2,5)$ existieren die 25 Geraden $y = mx + b$ \quad ($m, b \in \Zm{5}$).
Welche 5 Geraden fehlen in dieser Aufzählung?
\end{enumerate}
\end{aufgabe}

\begin{loesung}
\begin{enumerate}[label=\alph*)]
\item $y = 3x + 2$ : $(0, 2), (1, 0), (2, 3), (3, 1), (4, 4)$
% KORRIGIERT: Rechnung in Z_5 und y-Wert nicht angegeben
\item In $\Zm{5}$: $3x + 2 = 4x + 1 \Rightarrow \quad x - 1 = 0 \Rightarrow \quad x = 1,\ y = 3 \cdot 1 + 2 = 5 \equiv 0 \pmod 5 \Rightarrow \quad S = (1, 0)$
\item $x = k$ mit den Punkten $(k, 0)$, $(k, 1)$, $(k, 2)$, $(k, 3)$, $(k, 4)$ \quad ($k \in \Zm{5}$)
\end{enumerate}
\end{loesung}

\subsection{Bijektionen}

\begin{aufgabe}[Aufgabe 4]
Betrachten ein $4 \times 4$ Gitter.
% GRAFIK: Quadratisches Gitter aus 4 x 4 Kästchen (5 x 5 Gitterpunkte). Rot eingezeichnet ist ein Weg von der Ecke links unten zur Ecke rechts oben: r, o, r, o, o, r, r, o (r = eine Einheit nach rechts, o = eine Einheit nach oben).

Wie viele Wege entlang der Gitterlinien gibt es, die die Ecke links unten mit der Ecke rechts oben verbinden? Erlaubt sind nur Teilstrecken von links nach rechts oder von unten nach oben.
\end{aufgabe}

\begin{loesung}
Diese Wege setzen sich aus insgesamt 8 Strecken zusammen, wobei jeweils 4 Strecken nach oben ($o$) und nach rechts ($r$) verlaufen.

Der rot eingezeichnete kürzeste Weg kann geschrieben werden als
\[
r\ o\ r\ o\ o\ r\ r\ o
\]
Man hat also eine Bijektion zwischen den Wegen und der Menge der Kombinationen von 4 Strecken aus 8. (Man wählt die 4 Plätze auf denen $r$ steht.)
\[
\Rightarrow\quad \binom{8}{4} = \frac{8 \cdot 7 \cdot 6 \cdot 5}{1 \cdot 2 \cdot 3 \cdot 4} = 7 \cdot 2 \cdot 5 = 70
\]
Das sind keine Kombinationen mit Wiederholung, da bei Kombinationen die Anordnung keine Rolle spielt.
\end{loesung}

\begin{aufgabe}[Aufgabe 5]
Wie viele solche Wege gibt es für ein $m \times n$ Gitter?
\end{aufgabe}

\begin{loesung}
Die Wege setzen sich aus insgesamt $m+n$ Strecken zusammen, wobei jeweils $m$ Strecken nach oben und jeweils $n$ nach rechts verlaufen. Man hat also aus $m + n$ Strecken $m$ nach oben zu wählen.
\[
\binom{m+n}{n} = \binom{m+n}{m}
\]
\end{loesung}

\subsection{Binomischer Lehrsatz}

\begin{aufgabe}[Aufgabe 6]
% KORRIGIERT: "für natürliche Zahlen k" ohne k >= 1 (für k = 0 falsch)
Berechnen Sie für natürliche Zahlen $k \geq 1$ die Summe $\binom{k}{1} - \binom{k}{2} + \cdots + (-1)^{k-1}\binom{k}{k}$.
\end{aufgabe}

\begin{loesung}
Setzen im binomischen Lehrsatz $(a + b)^k = \ldots$ \quad $a = 1$, \ $b = -1$
% KORRIGIERT: "- + ..." und fehlende Begründung für "= 0"
\begin{align*}
\Rightarrow\quad & 1 - \binom{k}{1} + \binom{k}{2} - \cdots + (-1)^k \binom{k}{k} = (1 - 1)^k = 0 \quad (\text{da } k \geq 1)\\
\Rightarrow\quad & \binom{k}{1} - \binom{k}{2} + \cdots + (-1)^{k-1} \binom{k}{k} = 1 \qquad \square
\end{align*}
\end{loesung}

\begin{aufgabe}[Aufgabe 7]
Beweisen Sie durch Modifikation der Identität in Folgerung c)
\[
\binom{m}{0}\binom{n}{r} + \binom{m}{1}\binom{n}{r-1} + \cdots + \binom{m}{r}\binom{n}{0} = \binom{m+n}{r} \qquad r > 0,\ m \geq r,\ n \geq r
\]
\end{aufgabe}

\begin{loesung}
Betrachten den Koeffizienten von $x^r$ in der Identität
\[
(1 + x)^m (1 + x)^n = (1 + x)^{m+n}
\]
Die linke Seite ist das Produkt zweier Faktoren der Form
\[
1 + \binom{m}{1}x + \cdots + \binom{m}{k}x^k + \cdots + x^m, \quad 1 + \binom{n}{1}x + \cdots + \binom{n}{l}x^l + \cdots + x^n
\]
Die Potenz $x^r$ kommt zustande, wenn der Term $\binom{m}{k}x^k$ mit $\binom{n}{r-k}x^{r-k}$ multipliziert wird, $0 \leq k \leq r$, also
\[
\binom{m}{0}\binom{n}{r} + \cdots + \binom{m}{k}\binom{n}{r-k} + \cdots + \binom{m}{r}\binom{n}{0} = \binom{m+n}{r}. \qquad \square
\]
\end{loesung}
